\documentclass[11pt]{article}
\usepackage[a4paper,margin=25mm]{geometry}
\usepackage{amsmath,amssymb,amsthm,mathtools,hyperref,microtype}
\hypersetup{colorlinks=true,urlcolor=blue,linkcolor=blue}
\newtheorem{theorem}{Theorem}
\newtheorem{proposition}[theorem]{Proposition}
\newtheorem{corollary}[theorem]{Corollary}
\newcommand{\Ric}{\operatorname{Ric}}
\newcommand{\Rm}{\operatorname{Rm}}
\newcommand{\Conv}{\operatorname{Conv}}
\newcommand{\CP}{\mathbb{CP}}
\newcommand{\R}{\mathbb R}
\newcommand{\orb}{\mathrm{orb}}
\title{Einstein Defect at the Orbifold Boundary:\\Multiple-Core Obstructions and Actual Subthreshold Degeneration}
\author{Research continuation}
\date{9 October 2026}
\begin{document}
\maketitle
\begin{abstract}
We continue the smooth-attainment question for the four-dimensional functional
$P(g)=\int|\Ric_g^\circ|^2$. Actual monopole supply transfers to every metric
on a fixed orbifold obtained by replacing finitely many Ricci-flat ALE cores.
The entire supplied-class span has anti-self-dual restrictions on each core;
this preserves the convex-hull budget and yields an orientation-sensitive
curvature obstruction. Conversely, ordinary-orientation rational balls really
embed in every single blow-up. Combining a classical Godeaux degeneration with
the project's negative K\"ahler--Einstein Burns insertion gives a fixed
non-Einstein closed manifold with bubbling sequences strictly below
$8\pi^2$. This separates a strict upper bound from smooth attainment.
No global minimum or resolution of the four outstanding questions is asserted.
\end{abstract}

\section{Definitions and inputs}
All manifolds are oriented. Curvature norms are fixed by
\begin{align}
8\pi^2\chi(M)&=\int_M\left(|W^+|^2+|W^-|^2+\frac{s^2}{24}
                           -\frac{|\Ric^\circ|^2}{2}\right),\label{eq:gb}\\
12\pi^2\tau(M)&=\int_M(|W^+|^2-|W^-|^2).
\end{align}
Write $T(M)=2\chi(M)+3\tau(M)$ and
$P_*(M)=\inf_g P(g)$, with the infimum over all smooth metrics.
For an orbifold, $\chi_{\rm curv}$ and $\tau_{\rm curv}$ denote the
corresponding curvature integrals. In particular $\tau_{\rm curv}$ is
\emph{not} necessarily the signature of its real intersection form.

The inherited analytic input for Section~\ref{sec:subthreshold} is the
project's strict Burns insertion on a negative K\"ahler--Einstein background.
Its statement is in the supplied \emph{01\_Research\_Synthesis.tex},
Section~3, equations (insertion), (phi), and (kernel): in a kernel direction of
$\mathsf R^+=W^++sI/12$,
\begin{equation}
 P(g_\varepsilon)=P(g)+8\pi^2
 -4\pi^2|T_p|^2\varepsilon^2\log(1/\varepsilon)+O(\varepsilon^2).
 \label{eq:burns}
\end{equation}
Here $T_p$ is the trace-free restriction of $\mathsf R^+$ to the orthogonal
two-plane. On a K\"ahler--Einstein surface, the eigenvalues of $\mathsf R^+$
are $(s/4,0,0)$, so $|T_p|^2=s^2/32>0$ when $s<0$.
We use this project result as an input, not as a newly verified literature theorem.

\section{Actual monopole supply through several ALE cores}
Let $\mathcal C$ be a finite collection of determinant classes of fixed
spin$^c$ structures on $M$, each admitting an unperturbed Seiberg--Witten
solution for every smooth metric. Suppose compact cores of finitely many
oriented Ricci-flat ALE spaces $(N_i,g_i)$ embed disjointly in $M$.
Their ends are $S^3/\Gamma_i$, with positive ALE decay order and two
differentiated estimates. Replacing the cores by their marked quotient cones
gives a compact orbifold $O$ with isolated quotient singularities.
We use compatible orbifold spin$^c$ extensions of the boundary structures;
the finite isotropy representations are retained throughout.

Because the links are real homology spheres, there is an orthogonal splitting
\[
H^2(M;\R)=H^2_{\orb}(O;\R)\oplus\bigoplus_i H^2(N_i;\R).
\]
Write $v=(v_0;v_1,\ldots,v_r)$ and set
\[
B_O=\max_{v\in\Conv(\mathcal C)}v_0^2,
\qquad D_i=\frac{1}{2\pi^2}\int_{N_i}|W^+|^2.
\]

\begin{theorem}\label{thm:transfer}
Assume $B_O>0$. The restriction to each $N_i$ of the entire real span of
$\mathcal C$ is represented by $L^2$ anti-self-dual harmonic two-forms.
Consequently
\begin{equation}
 B_O\ge\max_{v\in\Conv(\mathcal C)}v^2.
 \label{eq:convex}
\end{equation}
Every smooth orbifold metric $g$ on $O$ satisfies
\begin{equation}
\boxed{P_{\orb}(g)\ge8\pi^2\left(\frac23 B_O-T(M)+\sum_iD_i\right).}
\label{eq:main}
\end{equation}
\end{theorem}

\begin{proof}
Fix an arbitrary orbifold metric $g$. Near each quotient point its lift has
$g=\delta+O(r^2)$ in normal coordinates, with two differentiated bounds.
If the ALE order is $\nu>0$, its rescaling by $a^2$ has error
$O(a^\nu r^{-\nu})$. Choose $r_a=a^q$ with
$0<q\le\nu/(\nu+2)$ and use a radial cutoff on
$r_a<r<2r_a$. The annular curvature satisfies
\[
 |\Rm|\le C\left[1+a^\nu r_a^{-\nu-2}
                 +(r_a+a^\nu r_a^{-\nu-1})^2\right]\le C.
\]
The resulting smooth metrics $g_a$ on $M$ have uniformly bounded volume,
a uniform lower scalar-curvature bound, and vanishing annular curvature
energy. The exact inner metrics are Ricci-flat, hence
\begin{equation}
 P(g_a)\longrightarrow P_{\orb}(g).
 \label{eq:smoothing}
\end{equation}
These are direct metric interpolations, not Einstein desingularizations.

Choose actual solutions for every supplied class. The maximum principle
bounds the spinors uniformly, hence bounds $F_A^+$ pointwise. The identity
\[
 \int_M(|F_A^+|^2-|F_A^-|^2)=4\pi^2a^2
\]
then bounds the determinant curvature in $L^2$. On compact subsets of
$O_{\rm reg}$, abelian gauge compactness and the equations give smooth
subsequential limits. On each covering ball an isolated point is removable:
finite-energy abelian curvature permits a $W^{1,2}$ Coulomb gauge, and the
bounded spinor and the elliptic equations bootstrap the fields across the
point. The extension respects the finite deck action; trivial holonomy is
not assumed. The limiting real classes are the restrictions $a_0$.
Since $g$ was arbitrary, this supplies actual fields for every orbifold
metric, including the conformal representatives used below.

On a rescaled core the same two-form satisfies
$|F_A^+|_{g_i}\le Ca^2\to0$. Its scale-invariant $L^2$ bound gives a
weak limit $F_i$ with $dF_i=0$ and $F_i^+=0$.
Elliptic regularity makes it an $L^2$ harmonic form.
Pairing against compactly supported closed two-forms identifies its class
as $a_i$. ALE $L^2$ Hodge theory identifies this cohomology with
$\operatorname{im}(H_c^2(N_i)\to H^2(N_i))$, which here equals $H^2(N_i)$.
In particular $a_i^2\le0$. A finite diagonal extraction places the limits
of all supplied classes on the same fixed core metrics. Their real linear
combinations are also anti-self-dual, so $v_i^2\le0$ for every
$v\in\operatorname{span}_{\R}\mathcal C$. This proves \eqref{eq:convex}.

For each core define
\[
 e_i=\frac1{8\pi^2}\int_{N_i}(|W^+|^2+|W^-|^2),\qquad
 t_i=\frac1{12\pi^2}\int_{N_i}(|W^+|^2-|W^-|^2).
\]
Passing to the limit in the closed characteristic identities along $g_a$
gives $\chi_{\rm curv}(O)=\chi(M)-\sum e_i$ and
$\tau_{\rm curv}(O)=\tau(M)-\sum t_i$.
Thus
\begin{equation}
 T_{\rm curv}(O)=T(M)-\sum_i(2e_i+3t_i)=T(M)-\sum_iD_i.
 \label{eq:charge}
\end{equation}

For any fixed $g$, convexity of the norm on $H_g^+$ implies that some
supplied class has self-dual square at least $B_O$.
The classical mixed monopole inequality, applied in invariant orbifold
charts, is
\[
 \int_O(s-\sqrt6|W^+|)^2\ge72\pi^2B_O.
\]
Weighted Cauchy--Schwarz gives
\[
 \int_O\left(\frac{s^2}{24}+2|W^+|^2\right)
 \ge\frac{8\pi^2}{3}B_O.
\]
Finally $4\pi^2T_{\rm curv}(O)$ equals the last integral minus
$P_{\orb}/2$. Combining with \eqref{eq:charge} proves \eqref{eq:main}.
\end{proof}

The ingredients are classical curvature inequalities, abelian compactness,
removability and ALE Hodge theory. The content assembled here is the
full-span restriction, the preserved convex budget and the multiple-core
obstruction. No classification of all Ricci-flat ALE spaces is used.

\begin{corollary}
Suppose $M=X\#k\overline{\CP}^{\,2}$ has all supplied sign classes
$K+\sum_{j=1}^k\epsilon_jE_j$, where $K^2=A>0$,
$K\cdot E_j=0$, $E_i\cdot E_j=-\delta_{ij}$ and $T(M)=A-k$.
Set $\Delta_K=-\sum_i K_i^2\ge0$. Then
\[
 P_{\orb}\ge8\pi^2\left(k-\frac A3+\sum_iD_i+\frac23\Delta_K\right).
\]
Consequently a sequence with $\limsup P<8\pi^2k$ cannot have this
smooth compact orbifold limit if
$\sum D_i+\frac23\Delta_K\ge A/3$.
\end{corollary}
Indeed the average sign class is $K$, and lower semicontinuity on compact
regular subsets gives $P_{\orb}\le\liminf P$. A neck energy equality is
unnecessary. For $r$ reversed Ricci-flat ASD rational balls with end groups
of orders $d_i^2$, the criterion becomes
\begin{equation}
 A\le12\sum_{i=1}^r(1-d_i^{-2});\qquad A\le9r\ \text{suffices}.
 \label{eq:reverse}
\end{equation}

\section{Ordinary rational balls are actual competing topologies}
\begin{proposition}
Every $X\#\overline{\CP}^{\,2}$ contains the ordinary oriented rational
ball $B_{n,1}$ for each odd $n\ge3$. Its fundamental group maps trivially
into that of the ambient manifold. It carries a Ricci-flat K\"ahler ALE
metric in the indicated orientation with
\[
 P(N)=0,\qquad W^+=0,\qquad
 \int_N|W^-|^2=8\pi^2(1-n^{-2}).
\]
\end{proposition}
\begin{proof}
Reversing the orientation of a nonsingular conic in $\CP^2$ gives a sphere
of square $-4$ in the $\overline{\CP}^{\,2}$ summand, disjoint from the
connected-sum ball. Khodorovskiy's theorem embeds $B_{n,1}$ into its
normal disk bundle $V_{-4}$ for odd $n\ge3$. This bundle is simply
connected, proving the claim about fundamental groups. Suvaina's theorem
supplies the Ricci-flat K\"ahler ALE metric with end group of order $n^2$.
Scalar-flat K\"ahler geometry gives $W^+=0$ and ALE Gauss--Bonnet gives
the displayed energy.
\end{proof}
This combines classical theorems; it is not a new rational-ball embedding.
The smoothing \eqref{eq:smoothing} produces actual fixed-manifold
degenerations through each such core. Both correction terms in the
corollary vanish: $D_i=0$ and $H^2(N_i;\R)=0$.
Thus ordinary ASD rational balls are invisible to these two improvements,
whereas reversed ones satisfy \eqref{eq:reverse}.

\section{An actual subthreshold bubbling sequence}\label{sec:subthreshold}
\begin{proposition}[Using the inherited strict insertion theorem]
There is a numerical Godeaux surface $Y$ with $H_1(Y;\mathbb Z)=\mathbb Z_4$
such that $M=Y\#\overline{\CP}^{\,2}$ admits smooth metrics $g_a$ and a
constant $q_0$ satisfying
\[
 \frac{16\pi^2}{3}\le P_*(M)\le q_0<8\pi^2,
 \qquad P(g_a)\longrightarrow q_0,
\]
while a nonflat ordinary Ricci-flat rational-ball bubble forms.
The manifold $M$ admits no Einstein metric.
\end{proposition}
\begin{proof}
Kazanova's Proposition 2.6 constructs a $\mathbb Q$-Gorenstein degeneration
of such a smooth $Y$ to a surface $Z$ with ample canonical divisor and
exactly one Wahl singularity $\frac14(1,1)$. Its local smoothing is the
ordinary $B_{2,1}$. Orbifold Aubin--Yau gives a negative K\"ahler--Einstein
metric $h_0$ on $Z$. At a different, smooth point apply \eqref{eq:burns}
inside a fixed ordinary chart. The negative logarithmic term gives a
metric $h$ on the one-point blow-up orbifold $O$ with
$q_0=P_{\orb}(h)<8\pi^2$. It equals $h_0$ near the Wahl point.
The insertion is local, so its proof is unaffected by the distant quotient
point.

Insert the scale-$a$ Ricci-flat ALE smoothing of $B_{2,1}$ at that point.
Its order-four expansion and the bulk normal-coordinate expansion are
\[
 a^2g_N=\delta+O(a^4r^{-4}),\qquad h=\delta+O(r^2).
\]
Matching at $\rho=a^{1/2}$ gives annular curvature $O(1)$ and annular
energy $o(1)$. The exact inner region has $P=0$; the removed bulk ball
also has $P=0$. Hence $P(g_a)\to q_0$. The smoothing topology and the
blow-up away from the singular point give the same fixed smooth manifold
$M$ for all small $a>0$. After rescaling, the inserted region exhausts
the nonflat ALE metric, so curvature is unbounded in the fixed bulk scale.

Here $K_Y^2=A=1$ and $T(M)=0$. The classical all-metric monopole bound
is $P(g)\ge8\pi^2(\frac23A-T)=16\pi^2/3$, excluding Einstein metrics.
This bound and the constructed sequence prove the claimed interval.
\end{proof}

The example is a torsion-$\mathbb Z_4$ Godeaux surface, not the classical
torsion-$\mathbb Z_5$ surface. Kazanova's Corollary 2.5 excludes this
particular ample-canonical single-Wahl degeneration for the latter family.
That algebraic exclusion alone says nothing about all smooth embeddings
or all Riemannian degenerations on it.

\section{The remaining comparison is variational}
The preceding proposition shows that the sublevel $P<8\pi^2$ need not be
smoothly compact even on a fixed non-Einstein example. It does not say that
$q_0=P_*(M)$, so it is neither a nonattainment theorem nor a solution of Q3.

For a specified orbifold boundary obtained by these Ricci-flat smoothings,
\eqref{eq:smoothing} gives
\[
 P_*(M)\le P_{*,\orb}(O).
\]
If a genuine minimizing sequence has this compact orbifold limit, lower
semicontinuity gives the reverse comparison through its limit metric.
Thus that branch requires equality with an optimized singular problem.
To exclude it one needs a strict variational comparison, or information
about genuine minimizing sequences that is absent from an arbitrary
sublevel bound. Partial collapse remains a separate possible branch.

The multiple-core theorem strengthens the range already excluded. The
subthreshold construction proves that the complementary orientation is a
real possibility at the relevant energy scale. These are complementary
conclusions about the same attainment problem, not two new research topics.

\begin{thebibliography}{9}
\bibitem{LeBrun} C. LeBrun, \emph{Ricci Curvature, Minimal Volumes, and
Seiberg--Witten Theory}, \href{https://arxiv.org/abs/math/0003068}{arXiv:math/0003068};
\emph{Four-Manifolds, Curvature Bounds, and Convex Geometry},
\href{https://arxiv.org/abs/math/0611450}{arXiv:math/0611450}.
\bibitem{HHM} T. Hausel, E. Hunsicker and R. Mazzeo,
\emph{Hodge cohomology of gravitational instantons}, Duke Math. J. 122
(2004), 485--548, \href{https://arxiv.org/abs/math/0207169}{arXiv:math/0207169}.
\bibitem{Khod} T. Khodorovskiy, \emph{Smooth Embeddings of Rational Homology
Balls}, Topology Appl. 161 (2014), 386--396, Theorems 1.2--1.3,
\href{https://arxiv.org/abs/1212.4391}{arXiv:1212.4391}.
\bibitem{Owens} B. Owens, \emph{Equivariant embeddings of rational homology
balls}, Theorem 2, \href{https://www.maths.gla.ac.uk/~bowens/papers/balls.pdf}{author manuscript}.
\bibitem{Suvaina1} I. Suvaina, \emph{ALE Ricci-flat K\"ahler metrics and
deformations of quotient surface singularities}, Ann. Global Anal. Geom. 41
(2012), 109--123, Proposition 3.3,
\href{https://arxiv.org/abs/1102.2672}{arXiv:1102.2672}.
\bibitem{Suvaina2} I. Suvaina, \emph{The Yamabe invariant of a class of
symplectic 4-manifolds}, Proposition 2.4 and Section 3,
\href{https://arxiv.org/abs/1410.1236}{arXiv:1410.1236}.
\bibitem{Wright} E. P. Wright, \emph{Quotients of gravitational instantons},
Theorem A and Section 8, \href{https://arxiv.org/abs/1102.2442}{arXiv:1102.2442}.
\bibitem{Kazanova} A. Kazanova, \emph{Degenerations of Godeaux surfaces and
exceptional vector bundles}, Proposition 2.6 and Corollary 2.5,
\href{https://arxiv.org/abs/1402.0254}{arXiv:1402.0254}.
\end{thebibliography}
\end{document}
